\documentclass[a4paper,12pt,oneside]{report}
\usepackage[utf8]{inputenc}
\usepackage[english]{babel}
\usepackage{geometry}
\usepackage{amsmath}
\geometry{a4paper, margin=1in}
\usepackage{tocloft}
\usepackage{graphicx}
\usepackage{setspace}
\usepackage{xcolor}
\usepackage{float}
\usepackage{hyperref}
\usepackage{caption}
\usepackage{longtable}
\usepackage{booktabs}
\usepackage{array}
\usepackage{listings}
\usepackage{cite}
\usepackage{ragged2e}
\usepackage{pdfpages}
\usepackage{indentfirst}
\usepackage{etoolbox}
\usepackage{fancyhdr}

\definecolor{bitblue}{RGB}{0, 70, 127}
\definecolor{codegray}{rgb}{0.95,0.95,0.95}

\lstset{
  backgroundcolor=\color{codegray},
  basicstyle=\ttfamily\small,
  breaklines=true,
  frame=single,
  language=Python
}

\hypersetup{
    colorlinks=true,
    linkcolor=black,
    citecolor=blue,
    urlcolor=cyan
}

\onehalfspacing

\setlength{\headheight}{13.6pt}
\fancyhf{}
\pagestyle{fancy}
\fancyhead[L]{\small\textit{Use of AI for the Analysis of Facial Expressions}}
\fancyhead[R]{\small\textit{of Internally Displaced Children}}
\renewcommand{\headrulewidth}{0.4pt}
\fancyfoot[L]{\small\textbf{SILGA Patricia} \quad Promotion: 2023--2026}
\fancyfoot[R]{\small Page \thepage}
\renewcommand{\footrulewidth}{0.4pt}

\renewcommand{\thepage}{\roman{page}}
\setcounter{tocdepth}{2}

\renewcommand{\cfttoctitlefont}{\hfil}
\renewcommand{\cftloftitlefont}{\hfil}
\renewcommand{\cftlottitlefont}{\hfil}
\renewcommand{\cftbeforechapskip}{10pt}
\renewcommand{\cftbeforesecskip}{5pt}
\renewcommand{\cftbeforesubsecskip}{3pt}

\begin{document}

%% ============================================================
%% COVER PAGE
%% ============================================================
\includepdf[pages=1]{cover.pdf}

%% ============================================================
%% QUOTES
%% ============================================================
\section*{Quotes}
\addcontentsline{toc}{section}{Quotes}
\vspace*{3cm}
\begin{flushright}
\textit{``The good news is that the moment you decide that what you know is}\\
\textit{more important than what you have been taught to believe,}\\
\textit{you will have shifted gears in your quest for abundance.''}\\[0.5cm]
--- \textbf{Earle Nightingale}\\[2cm]
\textit{``Artificial Intelligence is the new electricity.''}\\[0.5cm]
--- \textbf{Andrew Ng}
\end{flushright}
\clearpage

%% ============================================================
%% DEDICATION
%% ============================================================
\section*{Dedication}
\addcontentsline{toc}{section}{Dedication}
\vspace*{2cm}
\begin{flushright}
\textit{To my beloved brother, \textbf{SILGA Pengdwende Fabrice} ---}\\[0.3cm]
\textit{my pillar, my greatest support, and my most faithful ally}\\
\textit{in every battle of life. This work is as much yours as it is mine.}\\[1cm]
\textit{To my dear parents, whose love, sacrifice, and endless patience}\\
\textit{have been the silent force behind every step of this journey.}\\
\textit{With all my gratitude, my love, and my deepest respect.}\\[1cm]
\textit{To all my supervisors and teachers who guided me}\\
\textit{from my very first year to this final milestone ---}\\
\textit{thank you for your knowledge, your patience,}\\
\textit{and your dedication to shaping who I am today.}\\[1.5cm]
\textbf{SILGA Patricia}
\end{flushright}
\clearpage

%% ============================================================
%% ACKNOWLEDGEMENTS
%% ============================================================
\section*{Acknowledgements}
\addcontentsline{toc}{section}{Acknowledgements}
\vspace{0.5cm}
\justifying

I would like to express my sincere gratitude to all those who contributed, in one
way or another, to the successful completion of this thesis.

First and foremost, I thank the Almighty God for the gift of life, health, and
strength throughout this academic journey.

My deepest gratitude goes to my thesis supervisor, \textbf{Mr. Roland TOUGMA},
for his constant availability, valuable guidance, and constructive feedback
throughout this work. His patience and dedication have been truly inspiring.

I would like to express my heartfelt gratitude to \textbf{Dr. NIKIEMA Lebian
Wilfried}, whose wholehearted support and dedication have been invaluable
throughout this project. His encouragements and technical guidance have
significantly contributed to the realization of this work.

To the CEO of BIT, \textbf{Ms. Charlotte Rahe}, and the General Manager,
\textbf{Mr. Baslayi Tindano}, I express my profound gratitude for their
commitment to providing high-quality education and for creating an environment
that fosters academic excellence and innovation.

To all the faculty members of BIT, thank you for your dedication and commitment
throughout these three years.

To my beloved brother, \textbf{SILGA Pengdwende Fabrice}, and my dear parents ---
your love, sacrifices, and unconditional support have been my greatest source of
motivation. This achievement belongs to you as much as it does to me.

Finally, I warmly thank all my friends, classmates, and every person who supported
me morally and intellectually throughout this journey.
\clearpage

%% ============================================================
%% PREFACE
%% ============================================================
\section*{Preface}
\addcontentsline{toc}{section}{Preface}
\justifying

Burkina Institute of Technology (BIT) was inaugurated on October 15, 2018,
in Koudougou. It is a private institute founded by a German woman named
Susanne PERTL. The vision of BIT is to educate a new generation of leaders.
Students are trained with technical and entrepreneurial knowledge to help them
achieve this goal. That is also why the classes at BIT are taught in English,
to facilitate the internationalization of students.\\

BIT offers classes in three majors for the Bachelor's degree and one major for
the Master's degree. The majors for the Bachelor's Degree are:\\

\textbf{$\bullet$ Computer Science:} This was the first major when the institute
was inaugurated. There are courses in web, mobile, and desktop development, as
well as basics in networking, artificial intelligence, and data analysis.\\

\textbf{$\bullet$ Electrical Engineering:} This was the second major, opened in
2019 with a specialization in renewable energies. All the necessary courses
related to renewable energy in the electrical engineering field are provided.\\

\textbf{$\bullet$ Mechanical Engineering:} Introduced in 2021, with an opportunity
to specialize in either mining or agriculture during the final year.\\

BIT also has a laboratory where students carry out practical experiments. All of
these programs are completed in three (3) years and are open to anyone who holds
a scientific baccalaureate diploma.\\

The Master's Degree major, introduced in 2023, is focused on Artificial
Intelligence. It provides strong skills to help students master this promising
field.\\

At the end of the 3-year program, a thesis defense is required to obtain the
Bachelor's Degree. It is in this context that we undertook the realization of
this project on the use of Artificial Intelligence for the analysis of facial
expressions and biometric data of internally displaced children --- a project
that reflects both the technical and entrepreneurial spirit that BIT instills
in its students.

\clearpage

%% ============================================================
%% ABSTRACT (English)
%% ============================================================
\section*{Abstract}
\addcontentsline{toc}{section}{Abstract}
\justifying

Internal displacement caused by armed conflict and insecurity has become a major
humanitarian crisis in Burkina Faso, where the number of internally displaced
persons exceeded two million by 2023. Internally displaced children constitute
one of the most vulnerable groups, often experiencing significant psychological
distress that remains largely undetected due to the lack of adequate screening
tools in the field. This thesis presents the design and implementation of an
artificial intelligence-based web application called EmotiAid, capable of
analyzing facial expressions to automatically detect the emotional states of
internally displaced children. The model was trained using the FER-2013 dataset,
covering seven basic emotions defined by Paul Ekman: happiness, sadness, anger,
fear, disgust, surprise, and neutrality. After preprocessing including face
detection via OpenCV Haar Cascade, normalization, and data augmentation, a
Convolutional Neural Network (CNN) achieved a validation accuracy of 67.53\%
on the FER-2013 benchmark, exceeding the academic standard of 65\%. The final
system includes a responsive web application built with React and FastAPI that
allows humanitarian field workers to capture facial images via camera or file
upload and receive instant emotional state predictions with confidence scores.
An automatic alert mechanism notifies agents when a child displays high fear or
sadness scores, recommending immediate psychosocial follow-up.

\vspace{0.5cm}
\textbf{Keywords:} Artificial Intelligence, Convolutional Neural Network,
Facial Expression Recognition, Internally Displaced Children, Emotion Detection,
EmotiAid, Psychosocial Support, Deep Learning.
\clearpage

%% ============================================================
%% RESUME (French)
%% ============================================================
\section*{R\'esum\'e}
\addcontentsline{toc}{section}{R\'esum\'e}
\justifying

Le d\'eplacement interne caus\'e par les conflits arm\'es et l'ins\'ecurit\'e
est devenu une crise humanitaire majeure au Burkina Faso, o\`u le nombre de
personnes d\'eplac\'ees internes a d\'epass\'e deux millions en 2023. Les
enfants d\'eplac\'es internes constituent l'un des groupes les plus vuln\'erables,
souffrant souvent de d\'etresses psychologiques importantes qui restent largement
ind\'etect\'ees, faute d'outils de d\'epistage adapt\'es sur le terrain. Ce
m\'emoire pr\'esente la conception et la mise en \oe{}uvre d'une application web
bas\'ee sur l'intelligence artificielle appel\'ee EmotiAid, capable d'analyser
les expressions faciales afin de d\'etecter automatiquement les \'etats
\'emotionnels des enfants d\'eplac\'es internes. Le mod\`ele a \'et\'e
entra\^in\'e \`a partir du dataset FER-2013, couvrant sept \'emotions de base
d\'efinies par Paul Ekman~: joie, tristesse, col\`ere, peur, d\'ego\^ut,
surprise et neutralit\'e. Apr\`es un pr\'etraitement rigoureux incluant la
d\'etection de visage via OpenCV Haar Cascade, la normalisation et l'augmentation
de donn\'ees, un r\'eseau de neurones convolutifs (CNN) a atteint une pr\'ecision
de validation de 67,53\% sur le benchmark FER-2013, d\'epassant le standard
acad\'emique de 65\%. Le syst\`eme final comprend une application web responsive
d\'evelopp\'ee avec React et FastAPI, permettant aux agents humanitaires de
capturer des images faciales via cam\'era ou t\'el\'echargement et de recevoir
imm\'ediatement une pr\'ediction de l'\'etat \'emotionnel avec un score de
confiance.

\vspace{0.5cm}
\noindent\textbf{Mots-cl\'es :} Intelligence Artificielle, R\'eseau de Neurones
Convolutifs, Reconnaissance d'Expressions Faciales, Enfants D\'eplac\'es
Internes, D\'etection d'\'Emotions, EmotiAid, Soutien Psychosocial,
Apprentissage Profond.
\clearpage

%% ============================================================
%% TABLE OF CONTENTS
%% ============================================================
\begingroup
\renewcommand{\contentsname}{\Huge Table of Contents}
\setlength{\parskip}{5pt}
\tableofcontents
\endgroup
\clearpage

%% ============================================================
%% LIST OF ACRONYMS
%% ============================================================
\section*{List of Acronyms and Abbreviations}
\addcontentsline{toc}{section}{List of Acronyms and Abbreviations}
\renewcommand{\arraystretch}{1.3}

\begin{longtable}{|p{3.5cm}|p{10.5cm}|}
\hline
\textbf{Acronym} & \textbf{Definition} \\
\hline
AI       & Artificial Intelligence \\
\hline
API      & Application Programming Interface \\
\hline
Axios    & Promise-based HTTP Client for JavaScript \\
\hline
BIT      & Burkina Institute of Technology \\
\hline
CNN      & Convolutional Neural Network \\
\hline
CORS     & Cross-Origin Resource Sharing \\
\hline
CSV      & Comma-Separated Values \\
\hline
DL       & Deep Learning \\
\hline
FastAPI  & Modern Python Web Framework for Building APIs \\
\hline
FER      & Facial Expression Recognition \\
\hline
GPU      & Graphics Processing Unit \\
\hline
HTTP     & HyperText Transfer Protocol \\
\hline
IDP      & Internally Displaced Person \\
\hline
JS       & JavaScript \\
\hline
JSON     & JavaScript Object Notation \\
\hline
JWT      & JSON Web Token \\
\hline
Keras    & High-Level Neural Networks API (TensorFlow) \\
\hline
ML       & Machine Learning \\
\hline
ORM      & Object-Relational Mapping \\
\hline
OpenCV   & Open Source Computer Vision Library \\
\hline
Pillow   & Python Imaging Library \\
\hline
Python   & High-Level General-Purpose Programming Language \\
\hline
React    & JavaScript Library for User Interfaces \\
\hline
ReLU     & Rectified Linear Unit \\
\hline
REST     & Representational State Transfer \\
\hline
RUP      & Rational Unified Process \\
\hline
SQL      & Structured Query Language \\
\hline
SQLite   & Lightweight Relational Database Management System \\
\hline
SQLAlchemy & Python SQL Toolkit and ORM \\
\hline
SVM      & Support Vector Machine \\
\hline
TailwindCSS & Utility-First CSS Framework \\
\hline
TensorFlow & Open Source Machine Learning Platform \\
\hline
UML      & Unified Modeling Language \\
\hline
UNHCR    & United Nations High Commissioner for Refugees \\
\hline
URL      & Uniform Resource Locator \\
\hline
Uvicorn  & Lightning-Fast ASGI Server for Python \\
\hline
Vite     & Frontend Build Tool \\
\hline
XAI      & Explainable Artificial Intelligence \\
\hline
2TUP     & Two-Track Unified Process \\
\hline
\end{longtable}
\clearpage

%% ============================================================
%% LIST OF FIGURES
%% ============================================================
\addcontentsline{toc}{section}{List of Figures}
\listoffigures
\clearpage

%% ============================================================
%% LIST OF TABLES
%% ============================================================
\addcontentsline{toc}{section}{List of Tables}
\listoftables
\clearpage

%% ============================================================
%% Switch to Arabic numerals
%% ============================================================
\renewcommand{\thepage}{\arabic{page}}
\setcounter{page}{1}

%% ============================================================
%% GENERAL INTRODUCTION
%% ============================================================
\newpage
\thispagestyle{fancy}
\addcontentsline{toc}{chapter}{General Introduction}
{\LARGE \bfseries \noindent General Introduction}
\vspace{2em}
\justifying

Since 2015, Burkina Faso has been severely affected by an escalating security
crisis driven by armed groups operating across its territory. This conflict has
triggered one of the most alarming internal displacement crises in West Africa.
According to the United Nations High Commissioner for Refugees (UNHCR, 2023),
the number of internally displaced persons in Burkina Faso exceeded two million,
with children representing more than half of this population. These children are
uprooted from their homes, separated from familiar environments, deprived of
schooling, and exposed to violence and loss.

The psychological consequences of forced displacement are profound. Displaced
children frequently exhibit symptoms of post-traumatic stress disorder,
depression, anxiety, and emotional dysregulation. However, the detection of these
psychological disturbances remains a major challenge in the field, primarily
because qualified mental health professionals are scarce in displacement camps,
and traditional assessment methods require structured clinical environments
unavailable in emergency humanitarian contexts.

Artificial Intelligence, and more specifically Deep Learning applied to facial
expression recognition, offers an innovative technological response to this
humanitarian challenge. Convolutional Neural Networks have demonstrated remarkable
performance in automatically detecting human emotions from facial images, achieving
accuracy rates exceeding 65\% on the international FER-2013 benchmark
\cite{goodfellow2013}. This thesis proposes the design and development of
EmotiAid, a web application enabling humanitarian agents to photograph an
internally displaced child and instantly obtain an AI-powered emotional state
analysis.

\section*{Problem Statement}
\addcontentsline{toc}{section}{Problem Statement}

The central research question of this thesis is: \textit{How can Artificial
Intelligence be used to automatically detect emotional states of internally
displaced children from facial photographs, in order to assist humanitarian
agents in identifying children requiring urgent psychological follow-up, in the
absence of specialized mental health professionals on the ground?}

\section*{Objectives}
\addcontentsline{toc}{section}{Objectives}

\textbf{General Objective:} To design, develop, and validate an artificial
intelligence system for automatic recognition of facial expressions of internally
displaced children, integrated into a web application accessible to humanitarian
agents in the field.

\noindent\textbf{Specific Objectives:}
\begin{itemize}
    \item \textbf{SO1:} Train a CNN on the FER-2013 dataset to recognize seven
    emotions with validation accuracy exceeding 65\%;
    \item \textbf{SO2:} Design and develop a secure RESTful API using FastAPI
    integrating the trained AI model with real-time inference capability;
    \item \textbf{SO3:} Develop an intuitive responsive web interface using React
    enabling agents to capture photographs, view results, and manage child profiles;
    \item \textbf{SO4:} Implement an automatic psychological alert mechanism
    triggered when fear or sadness scores exceed a critical threshold.
\end{itemize}

\section*{General Methodology}
\addcontentsline{toc}{section}{General Methodology}

This project adopts the 2TUP (Two-Track Unified Process) software development
methodology, structured around two parallel tracks converging towards a unified
solution. The AI component was developed using supervised deep learning on the
FER-2013 dataset comprising 35,887 grayscale images across seven emotion classes.
Training was performed on the Kaggle platform using a Tesla P100-PCIE-16GB GPU.
The web application follows a three-tier architecture: a React frontend, a FastAPI
backend, and an SQLite database.

\section*{Thesis Organization}
\addcontentsline{toc}{section}{Thesis Organization}

This thesis is organized into two main parts. \textbf{Part I} covers the
theoretical foundations in Chapter 1 and the methodology and conceptual approach
in Chapter 2. \textbf{Part II} presents the design and implementation in Chapter
3, and the results and analysis in Chapter 4. A general conclusion synthesizes
the contributions and outlines perspectives for future work.

\clearpage

%% ============================================================
%% PART I
%% ============================================================
\part*{Part I: Theoretical Foundations and State of the Art}
\addcontentsline{toc}{part}{Part I: Theoretical Foundations and State of the Art}
\addtocontents{toc}{\vspace{10pt}}

%% ============================================================
%% CHAPTER 1
%% ============================================================
\newpage
\thispagestyle{fancy}
\addcontentsline{toc}{chapter}{Chapter 1: General Information}
{\LARGE \bfseries \noindent Chapter 1: General Information}
\vspace{2em}
\renewcommand{\thesection}{\Roman{section}}
\justifying

\section{Context, Problematics and Objectives}

\subsection{Context}

The humanitarian situation in Burkina Faso has dramatically deteriorated since
2015. Children face unique vulnerabilities including heightened risk of physical
harm, psychological trauma, malnutrition, and loss of access to education.
Humanitarian organizations such as UNICEF and Save the Children have identified
early detection of psychological distress as a critical priority. However, the
scarcity of trained mental health professionals in the field, combined with
cultural and linguistic barriers, makes systematic psychological screening
nearly impossible with conventional methods.

\subsection{Problematics}

The following specific problems are identified:

\begin{itemize}
    \item \textbf{P1:} Mental health professionals are insufficient in number to
    serve the large displaced child population in the field;
    \item \textbf{P2:} Children in distress often cannot verbally express their
    psychological suffering, making clinical assessment difficult;
    \item \textbf{P3:} Existing AI-based facial expression systems are not adapted
    for deployment in low-resource humanitarian environments;
    \item \textbf{P4:} Most FER datasets are adult-centric and may not generalize
    well to child populations in African contexts.
\end{itemize}

\subsection{Objectives}

These problems lead to the following objectives: develop an AI model capable of
classifying seven basic emotions from facial images; build a simple web application
deployable with minimal infrastructure; enable non-specialist humanitarian agents
to conduct rapid emotional assessments; and contribute to the state of the art by
applying FER technology in a humanitarian context.

\section{General Information on Key Concepts}

\subsection{Facial Biometric Data}

Biometric data refers to measurable biological characteristics used to identify
or analyze individuals. Facial biometric data encompasses measurable features
extracted from the human face, including geometric measurements between facial
landmarks such as eyes, nose, and mouth, as well as texture patterns and dynamic
features such as micro-expressions. The face contains 43 muscles capable of
producing thousands of distinct expressions, mapped by the Facial Action Coding
System developed by Ekman and Friesen \cite{ekman1978} into discrete Action Units
forming the biological basis for emotion classification.

\subsection{Facial Expressions and Their Emotional Significance: Ekman's Model}

Paul Ekman's foundational research established that six basic emotions are
universally expressed through the same facial configurations regardless of
cultural background: happiness, sadness, anger, fear, disgust, and surprise
\cite{ekman1999}. A seventh category, neutrality, is added in computational
systems. These seven categories form the basis of the FER-2013 dataset and of
the EmotiAid classification system.

\begin{table}[H]
\centering
\caption{The Seven Basic Emotions and Their Facial Characteristics}
\begin{tabular}{|l|p{9cm}|}
\hline
\textbf{Emotion} & \textbf{Key Facial Indicators} \\
\hline
Happiness  & Raised lip corners, raised cheeks, crow's feet wrinkles \\
\hline
Sadness    & Inner brow raise, brow lowering, lip corner depression \\
\hline
Anger      & Brow lowering, upper eyelid raise, lip tightening \\
\hline
Fear       & Brow raise and pull together, upper eyelid raise, lip stretch \\
\hline
Disgust    & Nose wrinkling, upper lip raise, cheek raise \\
\hline
Surprise   & Brow raise, wide-open eyes, dropped jaw \\
\hline
Neutral    & No predominant muscular activation \\
\hline
\end{tabular}
\label{tab:emotions}
\end{table}

\subsection{Internally Displaced Children: Context and Psychological Vulnerabilities}

According to the UN Guiding Principles on Internal Displacement, internally
displaced persons are individuals forced to leave their homes due to armed
conflict, violence, or natural disasters, but who remain within their country's
borders. Children constitute a particularly vulnerable sub-group. The
psychological literature identifies specific risk factors including direct
exposure to violence, separation from caregivers, disruption of schooling, loss
of social networks, and chronic uncertainty about the future. These factors
create conditions conducive to post-traumatic stress disorder, depression, anxiety
disorders, and behavioral dysregulation.

\subsection{Fundamentals of Artificial Intelligence}

Artificial Intelligence refers to the simulation of human cognitive functions by
computational systems, encompassing capabilities such as learning, reasoning,
problem solving, and perception. The current wave of AI, driven by the
availability of large datasets, powerful GPUs, and advances in neural network
architectures, has produced systems that match or exceed human performance on
specific perceptual tasks \cite{lecun1998}.

\subsection{Deep Learning: A Key Sub-domain of AI}

Deep learning is a sub-field of machine learning using multi-layered artificial
neural networks to learn hierarchical representations of data. For image data,
Convolutional Neural Networks have emerged as the dominant architecture,
exploiting the spatial structure of images through convolutional filters detecting
local patterns such as edges, textures, and shapes.

\subsection{Presentation of the Proposed Solution: EmotiAid}

EmotiAid is a web-based application integrating a trained CNN model into a
humanitarian-oriented interface. The system allows a humanitarian field agent to
upload a photograph or capture an image via the device camera, receive an automatic
prediction of the child's dominant emotional state with confidence scores for all
seven classes, and trigger an automatic alert when fear or sadness exceeds 60\%.
Analysis results are saved to a database and accessible through a historical
dashboard and statistical report. The application is built on a FastAPI backend
and a React frontend, making it deployable on any standard web server and
accessible from any modern browser.

\subsection{Types of Emotions Detected by the Model}

The EmotiAid model detects the seven emotions of the Ekman model: Anger, Disgust,
Fear, Happiness, Neutrality, Sadness, and Surprise. Among these, Fear and Sadness
are classified as high-priority emotions triggering the psychological alert
mechanism, as they are most closely associated with post-traumatic stress and
acute psychological distress in the context of forced displacement.

\section{Ethical and Regulatory Considerations}

\subsection{Legal and Regulatory Framework}

The use of facial biometric data of children raises important legal and ethical
questions. The General Data Protection Regulation (GDPR) classifies biometric
data as sensitive personal data requiring explicit consent and specific security
measures. The UN Convention on the Rights of the Child establishes that any
technology used in relation to children must prioritize their best interests.

\subsection{Informed Consent and Data Protection}

In the context of EmotiAid, the humanitarian agent must obtain verbal consent
from the child's guardian before capturing and analyzing a facial image. Images
are processed in real time and are not retained in the database by default; only
the analysis results and associated metadata are stored. No image is shared with
third parties.

\subsection{Anonymization and Data Security}

All stored data is associated with an anonymized child identifier rather than a
full name. The application uses JWT authentication to ensure that only registered
agents can access analysis data.

\subsection{Ethical Risks and Limitations}

Several ethical risks are identified: misclassification leading to incorrect
psychological assessment; potential stigmatization of children based on detected
emotions; and data breach risks. These risks are mitigated by framing EmotiAid
as a support tool for trained agents rather than an autonomous diagnostic system.
The FER-2013 dataset, primarily composed of adult faces from Western contexts,
may introduce cultural and demographic biases that must be considered when
interpreting analysis results.

\section{State of the Art}

\subsection{Prior Work on AI and Emotion Recognition}

The field of automatic facial expression recognition has evolved significantly
over the past three decades. Early approaches relied on hand-crafted features
such as Local Binary Patterns (LBP) and Histogram of Oriented Gradients (HOG),
combined with classical classifiers such as Support Vector Machines. The
introduction of deep learning transformed the FER landscape. The FER-2013
challenge at ICML 2013, introduced by Goodfellow et al. \cite{goodfellow2013},
marked a turning point, establishing a large-scale benchmark comprising 35,887
labeled grayscale facial images across seven emotion classes.

\subsection{Existing Models in the Literature}

\textbf{FER+} \cite{barsoum2016} is an improved version of FER-2013, where each
image was annotated by ten annotators rather than one, producing more reliable
probabilistic labels and achieving higher model accuracy due to better annotation
quality.

\textbf{DeepFace} \cite{taigman2014}, developed by Facebook AI Research, uses a
nine-layer deep CNN trained on four million face images to achieve 97.35\%
accuracy on the Labeled Faces in the Wild benchmark. Although primarily designed
for identity verification, its architecture strongly influenced subsequent FER
systems.

\textbf{AffectNet} \cite{mollahosseini2017} is the largest publicly available
facial expression dataset with over one million images. Its ethnic and demographic
diversity makes it particularly relevant for systems targeting diverse populations.

\begin{table}[H]
\centering
\caption{Comparison of Key FER Models in the Literature}
\begin{tabular}{|l|l|l|l|}
\hline
\textbf{Model} & \textbf{Authors} & \textbf{Dataset} & \textbf{Accuracy} \\
\hline
FER-2013 baseline & Goodfellow et al. (2013) & FER-2013 & 65.0\% \\
\hline
FER+ & Barsoum et al. (2016) & FER+ & 81.20\% \\
\hline
DeepFace & Taigman et al. (2014) & LFW & 97.35\% \\
\hline
AffectNet & Mollahosseini et al. (2017) & AffectNet & 58.00\% \\
\hline
\textbf{EmotiAid CNN (This work)} & \textbf{SILGA (2026)} & \textbf{FER-2013} & \textbf{67.53\%} \\
\hline
\end{tabular}
\label{tab:models}
\end{table}

\subsection{Justification of the Proposed Solution}

The EmotiAid CNN achieves 67.53\% validation accuracy on FER-2013, surpassing the
65\% academic standard established by Goodfellow et al. \cite{goodfellow2013}.
While pre-trained models like VGG-Face and ResNet offer higher performance, they
require significant computational resources, making them unsuitable for CPU-based
deployment in humanitarian environments. The custom lightweight CNN architecture
developed in this project balances prediction accuracy with computational
efficiency, addressing a real-world deployment gap absent from the existing
literature.

\subsection*{Chapter Conclusion}

This chapter established the theoretical foundations of EmotiAid. The key concepts
underpinning the system were reviewed, including facial biometric data, Ekman's
universal emotion model, the vulnerability of internally displaced children, the
fundamentals of AI and CNNs, and the ethical framework governing the use of
biometric technology on vulnerable populations. The state of the art analysis
confirmed that a deployment gap exists in humanitarian contexts that EmotiAid
specifically addresses.

\clearpage

%% ============================================================
%% CHAPTER 2
%% ============================================================
\newpage
\thispagestyle{fancy}
\addcontentsline{toc}{chapter}{Chapter 2: Methodology and Conceptual Approach}
{\LARGE \bfseries \noindent Chapter 2: Methodology and Conceptual Approach}
\vspace{2em}
\renewcommand{\thesection}{\Roman{section}}
\justifying

\section{Methodology}

\subsection{Software Development Methodology: The 2TUP Model}

This project adopts the 2TUP (Two-Track Unified Process) methodology, an
adaptation of the Rational Unified Process specifically designed for
object-oriented software projects. The 2TUP model structures development into
two parallel tracks:

\begin{itemize}
    \item \textbf{The Functional Track:} Focuses on functional requirements, user
    needs, use cases, and business logic. For EmotiAid, this encompasses
    requirements analysis for the humanitarian agent interface and the emotion
    detection workflow definition.
    \item \textbf{The Technical Track:} Addresses technical architecture,
    frameworks, and infrastructure, including the CNN architecture selection, the
    React/FastAPI technology stack, and the REST API design.
\end{itemize}

The two tracks converge in the design and implementation phase, ensuring the final
system is both functionally complete and technically sound.

\subsection{Data Collection}

The primary dataset is \textbf{FER-2013} (Facial Expression Recognition 2013),
introduced at the ICML 2013 Challenges in Representation Learning workshop
\cite{goodfellow2013}.

\begin{table}[H]
\centering
\caption{FER-2013 Dataset Characteristics}
\begin{tabular}{|l|l|}
\hline
\textbf{Property} & \textbf{Value} \\
\hline
Total samples & 35,887 grayscale images \\
\hline
Image dimensions & 48 $\times$ 48 pixels \\
\hline
Number of classes & 7 (Angry, Disgust, Fear, Happy, Sad, Surprise, Neutral) \\
\hline
Training set & 28,709 images \\
\hline
Validation set & 3,589 images \\
\hline
Test set & 3,589 images \\
\hline
Format & Image files (JPG) \\
\hline
Source & Kaggle (msambare/fer2013) \\
\hline
\end{tabular}
\label{tab:fer2013}
\end{table}

\subsection{Data Preprocessing}

Raw FER-2013 data requires the following preprocessing steps:

\begin{enumerate}
    \item \textbf{Normalization:} Pixel values are scaled from [0, 255] to [0.0, 1.0].
    \item \textbf{Reshaping:} Arrays are reshaped to (48, 48, 1) for CNN input.
    \item \textbf{Label encoding:} Integer class labels are one-hot encoded.
    \item \textbf{Data augmentation:} Horizontal flipping, rotation ($\pm$10°),
    zoom (10\%), and brightness adjustment are applied to improve generalization.
\end{enumerate}

\subsection{Data Validation}

Before training, dataset integrity is verified through class distribution analysis
to identify imbalance, sample visualization to confirm label-image correspondence,
and statistical validation of normalized pixel value distributions.

\subsection{Model Development}

The CNN model is developed using TensorFlow 2.x and Keras. The development follows
an iterative process: a baseline architecture is first established and trained,
then progressively refined based on validation performance using EarlyStopping and
ReduceLROnPlateau callbacks.

\subsection{Model Training}

Training is performed on the Kaggle cloud platform using a Tesla P100-PCIE-16GB
GPU. The Adam optimizer with learning rate 0.001 is used, with categorical
cross-entropy as the loss function, a batch size of 64, and a maximum of 50 epochs.

\subsection{Evaluation Metrics}

The model is evaluated using accuracy, precision, recall, F1-score per class, and
a confusion matrix visualizing classification performance across all seven emotion
classes.

\section{Conceptual Approach}

\subsection{General System Architecture}

The EmotiAid system follows a three-tier client-server architecture: (1) the
Presentation Layer, a React.js responsive web application providing the user
interface; (2) the Application Layer, a FastAPI REST API handling authentication,
business logic, and AI model inference; and (3) the Data Layer, an SQLite database
storing user accounts, child profiles, and analysis histories.

\subsection{Functional Architecture}

The system supports the following main functional flows: authentication flow
(agent logs in, JWT token issued); analysis flow (image captured, sent to API,
preprocessed, CNN inference, result returned and saved); history flow (agent
queries past analyses); and statistics flow (aggregate emotion data displayed on
dashboard).

\subsection{Software Architecture}

The backend follows FastAPI's router-based modular architecture with separation
between routes, models, schemas, and database logic. The frontend uses React's
component-based architecture with functional components and hooks for state
management.

\subsection{Hardware Architecture}

Training was performed on Kaggle GPU cloud (Tesla P100-PCIE-16GB). Production
deployment requires a server with minimum 4-core CPU, 8 GB RAM, and 10 GB storage.
Client access requires only a mobile phone with a web browser and internet
connection.

\subsection{AI Inference Architecture}

During inference, the system executes the following pipeline: (1) image received
as a multipart form upload via HTTP POST; (2) decoded using Pillow and converted
to grayscale; (3) face detected and cropped using OpenCV Haar Cascade
\cite{viola2001}; (4) resized to 48$\times$48 and normalized to [0, 1]; (5) passed
through the loaded Keras CNN model; (6) Softmax output converted to emotion label
and confidence scores; (7) result serialized as JSON and returned to the frontend
with alert status.

\subsection{Technological Choice: Responsive Web Application}

The choice of a responsive web application was motivated by three factors:
(1) accessibility --- any device with a browser can access the system without
installation; (2) maintainability --- a single codebase serves all platforms;
(3) deployability --- the system can be hosted on a local server accessible via
local Wi-Fi without internet dependency.

\subsection*{Chapter Conclusion}

This chapter presented the 2TUP methodology and the conceptual architecture of
EmotiAid. The FER-2013 dataset provides a solid foundation for CNN training, and
the preprocessing pipeline ensures data quality. The three-tier architecture
establishes a clean separation between presentation, business logic, and data
layers.

\clearpage

%% ============================================================
%% PART II
%% ============================================================
\part*{Part II: Design and Implementation}
\addcontentsline{toc}{part}{Part II: Design and Implementation}
\addtocontents{toc}{\vspace{10pt}}

%% ============================================================
%% CHAPTER 3
%% ============================================================
\newpage
\thispagestyle{fancy}
\addcontentsline{toc}{chapter}{Chapter 3: Design, Modelling and Conceptual Approach}
{\LARGE \bfseries \noindent Chapter 3: Design, Modelling and Conceptual Approach}
\vspace{2em}
\renewcommand{\thesection}{\Roman{section}}
\justifying

\section{Functional and Technical Requirements}

\subsection{Functional Requirements}

\begin{table}[H]
\centering
\caption{Functional Requirements of EmotiAid}
\begin{tabular}{|l|p{7cm}|l|}
\hline
\textbf{ID} & \textbf{Requirement} & \textbf{Priority} \\
\hline
FR-01 & Agent authenticates with username and password & Must Have \\
\hline
FR-02 & Agent uploads a child photograph for analysis & Must Have \\
\hline
FR-03 & Agent captures an image via device camera & Must Have \\
\hline
FR-04 & System displays detected emotion and confidence scores & Must Have \\
\hline
FR-05 & System generates alert for fear or sadness $>$ 60\% & Must Have \\
\hline
FR-06 & Analysis results are automatically saved to database & Must Have \\
\hline
FR-07 & Agent consults analysis history & Should Have \\
\hline
FR-08 & Agent manages child profiles (create, view, history) & Should Have \\
\hline
FR-09 & Dashboard displays emotion distribution statistics & Should Have \\
\hline
FR-10 & Agent generates and prints a PDF statistical report & Could Have \\
\hline
\end{tabular}
\label{tab:requirements}
\end{table}

\subsection{Technical Requirements}

The system must satisfy the following technical constraints: emotion prediction
response time under 5 seconds; application accessible from any modern browser
without installation; responsive design for mobile and desktop; secure JWT-based
authentication; and REST API following standard HTTP conventions.

\section{UML Modelling}

\subsection{Use Case Diagram}

The system identifies three main actors. The \textbf{Humanitarian Agent} is the
primary user who logs in, captures or uploads images, views emotion analysis
results, manages child profiles, and consults the dashboard. The \textbf{AI
System} is an automated component that preprocesses images, runs CNN inference,
and returns predictions. The \textbf{Database} persists user accounts, child
profiles, and analysis records.

\begin{figure}[H]
\centering
\includegraphics[width=0.85\textwidth]{use_case_diagram.png}
\caption{Use Case Diagram of EmotiAid}
\label{fig:usecase}
\end{figure}

\subsection{Class Diagram}

The main entities of the system are:
\begin{itemize}
    \item \textbf{Agent:} id, username, password\_hash, nom\_complet, created\_at;
    \item \textbf{Child:} id, prenom, age\_approx, zone, agent\_id (FK), created\_at;
    \item \textbf{Analysis:} id, child\_id (FK), agent\_id (FK), emotion\_principale,
    scores (JSON), alerte (bool), created\_at;
    \item \textbf{EmotionModel:} model\_path, emotions, threshold, predict(),
    detectFace(), preprocess();
    \item \textbf{FaceDetector:} cascade, detect().
\end{itemize}

\begin{figure}[H]
\centering
\includegraphics[width=0.9\textwidth]{class_diagram.png}
\caption{Class Diagram of EmotiAid}
\label{fig:classdiagram}
\end{figure}

\subsection{Sequence Diagram}

The main sequence for an emotion analysis is: (1) agent opens the Analysis page;
(2) agent captures or uploads an image; (3) frontend sends HTTP POST to
\texttt{/predict} with the image file and JWT token; (4) FastAPI backend receives
the request, verifies the token, and calls the preprocessing pipeline; (5) OpenCV
detects and crops the face; (6) preprocessed image passed to the CNN model;
(7) model returns Softmax predictions; (8) backend formats JSON response with
emotion label, all scores, and alert status, saves to database; (9) frontend
renders the result card with emotion, confidence bars, and alert badge.

\begin{figure}[H]
\centering
\includegraphics[width=0.9\textwidth]{sequence_diagram.png}
\caption{Sequence Diagram of EmotiAid Analysis Flow}
\label{fig:sequence}
\end{figure}

\section{Development of the AI Model}

\subsection{Model Architecture}

The EmotiAid CNN is a custom architecture with 29 layers comprising four
convolutional blocks followed by a fully connected classifier.

\begin{table}[H]
\centering
\caption{EmotiAid CNN Architecture Summary}
\begin{tabular}{|l|l|l|}
\hline
\textbf{Block} & \textbf{Configuration} & \textbf{Output Shape} \\
\hline
Block 1 & Conv2D(32) + BN + MaxPool + Dropout(0.25) & 24$\times$24$\times$32 \\
\hline
Block 2 & Conv2D(64) + BN + MaxPool + Dropout(0.25) & 12$\times$12$\times$64 \\
\hline
Block 3 & Conv2D(128) + BN + MaxPool + Dropout(0.25) & 6$\times$6$\times$128 \\
\hline
Block 4 & Conv2D(256) + BN + MaxPool + Dropout(0.25) & 3$\times$3$\times$256 \\
\hline
Classifier & Flatten + Dense(512) + Dense(256) + Softmax(7) & 7 \\
\hline
\end{tabular}
\label{tab:cnn}
\end{table}

\subsection{Training Configuration}

\begin{itemize}
    \item Optimizer: Adam, learning rate 0.001;
    \item Loss function: Categorical cross-entropy;
    \item Batch size: 64; Maximum epochs: 50;
    \item Callbacks: ModelCheckpoint (save best model), EarlyStopping
    (patience=10), ReduceLROnPlateau (factor=0.5, patience=5);
    \item Training hardware: Kaggle GPU Tesla P100-PCIE-16GB.
\end{itemize}

\subsection{Training Process}

Training was completed in approximately 44 epochs before EarlyStopping was
triggered. The best model weights were saved as \texttt{emotion\_model\_best.h5}.
Training curves showed steady improvement with convergence observed after epoch 30.

\begin{figure}[H]
\centering
\includegraphics[width=0.8\textwidth]{training_curves.png}
\caption{Training and Validation Accuracy/Loss Curves}
\label{fig:training}
\end{figure}

\section{Development of the Application}

\subsection{Technologies Used}

\begin{table}[H]
\centering
\caption{Technology Stack for EmotiAid}
\begin{tabular}{|l|l|l|}
\hline
\textbf{Component} & \textbf{Technology} & \textbf{Version} \\
\hline
Backend framework & FastAPI & 0.104 \\
\hline
Authentication & Python-Jose + Passlib & JWT / bcrypt 4.0.1 \\
\hline
AI / Deep Learning & TensorFlow / Keras & 2.13 \\
\hline
Image processing & OpenCV, Pillow & 4.8 / 10.0 \\
\hline
Database ORM & SQLAlchemy + SQLite & 2.0 \\
\hline
Frontend framework & React.js + Vite & 18.2 \\
\hline
CSS framework & TailwindCSS & 3.0 \\
\hline
HTTP client & Axios & 1.4 \\
\hline
Language & Python / JavaScript & 3.12 / ES6+ \\
\hline
\end{tabular}
\label{tab:tech}
\end{table}

\subsection{Application Structure}

\begin{lstlisting}[caption={EmotiAid Project Structure}]
EmotiAid/
|-- Backend/
|   |-- main.py
|   |-- database.py
|   |-- models.py
|   |-- schemas.py
|   |-- routes/
|   |   |-- auth.py
|   |   |-- predict.py
|   |   |-- children.py
|   |   `-- analyses.py
|   `-- model/
|       `-- emotion_model_best.h5
`-- frontend/
    |-- src/
    |   |-- pages/
    |   |   |-- Login.jsx
    |   |   |-- Dashboard.jsx
    |   |   |-- Analyse.jsx
    |   |   |-- Historique.jsx
    |   |   |-- Rapport.jsx
    |   |   `-- ProfilEnfant.jsx
    |   `-- App.jsx
    `-- package.json
\end{lstlisting}

\subsection{Key Features of the Interface}

The EmotiAid interface provides six main pages. The \textbf{Login page} handles
agent authentication with JWT token storage. The \textbf{Dashboard} displays
key statistics and an emotion distribution chart. The \textbf{Analysis page}
offers camera capture and file upload with emotion results. The \textbf{History
page} lists all past analyses. The \textbf{Child Profiles page} manages
individual child records. The \textbf{Report page} generates a printable PDF
statistical report.

\subsection*{Chapter Conclusion}

This chapter presented the complete design of EmotiAid, from functional and
technical requirements through UML models to the CNN architecture and technology
stack. The following chapter presents the implementation results and their analysis.

\clearpage

%% ============================================================
%% CHAPTER 4
%% ============================================================
\newpage
\thispagestyle{fancy}
\addcontentsline{toc}{chapter}{Chapter 4: Implementation, Results and Analysis}
{\LARGE \bfseries \noindent Chapter 4: Implementation, Results and Analysis}
\vspace{2em}
\renewcommand{\thesection}{\Roman{section}}
\justifying

\section{Implementation}

\subsection{Backend Implementation}

The trained CNN model is loaded once at server startup and kept in memory to
minimize inference latency. The prediction route receives a multipart image
upload, runs the preprocessing pipeline (face detection via Haar Cascade,
grayscale conversion, resizing to 48$\times$48, normalization), and passes the
array to the model.

\begin{lstlisting}[caption={FastAPI Prediction Route (routes/predict.py)}]
from fastapi import APIRouter, File, UploadFile, Header
from tensorflow import keras
import numpy as np, cv2
from PIL import Image
import io

model = keras.models.load_model('model/emotion_model_best.h5')
EMOTIONS = ['angry','disgust','fear','happy','neutral','sad','surprise']

@router.post("/predict")
async def predict_emotion(file: UploadFile = File(...),
                          authorization: str = Header(None)):
    image_bytes = await file.read()
    face_bytes = detect_face(image_bytes)
    img_array = preprocess_image(face_bytes)
    predictions = model.predict(img_array, verbose=0)[0]
    emotion_index = int(np.argmax(predictions))
    scores = {EMOTIONS[i]: round(float(predictions[i])*100, 2)
              for i in range(len(EMOTIONS))}
    alerte = scores.get('fear',0) > 60 or scores.get('sad',0) > 60
    return {"emotion_principale": EMOTIONS[emotion_index],
            "scores": scores, "alerte": alerte}
\end{lstlisting}

\subsection{Frontend Implementation}

The React Analysis page implements both camera capture via the
\texttt{getUserMedia()} API and file upload, with automatic display of results
upon successful prediction.

\section{Results}

\subsection{Model Training Results}

The CNN model was trained on FER-2013 over 50 epochs with EarlyStopping. The best
validation accuracy of \textbf{67.53\%} was achieved at epoch 44.

\begin{table}[H]
\centering
\caption{EmotiAid CNN Training Results}
\begin{tabular}{|l|l|}
\hline
\textbf{Metric} & \textbf{Value} \\
\hline
Final Training Accuracy & 72.8\% \\
\hline
Best Validation Accuracy & \textbf{67.53\%} \\
\hline
Best Epoch & 44 \\
\hline
Model File Size & 21.9 MB \\
\hline
Total Parameters & $\sim$4.2 million \\
\hline
Training Duration & $\sim$45 minutes (Kaggle GPU) \\
\hline
\end{tabular}
\label{tab:results}
\end{table}

\subsection{Performance per Emotion Class}

\begin{table}[H]
\centering
\caption{Per-Class Performance on FER-2013 Test Set}
\begin{tabular}{|l|l|l|l|}
\hline
\textbf{Emotion} & \textbf{Precision} & \textbf{Recall} & \textbf{F1-Score} \\
\hline
Angry    & 0.61 & 0.60 & 0.61 \\
\hline
Disgust  & 0.46 & 0.46 & 0.46 \\
\hline
Fear     & 0.45 & 0.45 & 0.45 \\
\hline
Happy    & 0.88 & 0.88 & 0.88 \\
\hline
Neutral  & 0.63 & 0.63 & 0.63 \\
\hline
Sad      & 0.53 & 0.53 & 0.53 \\
\hline
Surprise & 0.77 & 0.77 & 0.77 \\
\hline
\textbf{Overall} & \textbf{0.67} & \textbf{0.67} & \textbf{0.67} \\
\hline
\end{tabular}
\label{tab:perclass}
\end{table}

\begin{figure}[H]
\centering
\includegraphics[width=0.8\textwidth]{confusion_matrix.png}
\caption{Confusion Matrix on FER-2013 Test Set}
\label{fig:confusion}
\end{figure}

\subsection{Web Application Results}

The EmotiAid web application was successfully developed and tested. Key
observations include: average prediction response time of 2.3 seconds on CPU;
image upload functional for JPEG and PNG formats; camera capture functional on
Chrome and Firefox; alert system correctly triggered for fear and sadness
predictions above 60\%; history and dashboard pages display data accurately.

\subsection{Real-World Prediction Test}

During functional testing, a real photograph was submitted through the EmotiAid
interface. The model returned: \textbf{Neutral 86.6\%}, Sad 7.6\%, Angry 3.6\%,
with no alert triggered. This result demonstrates the system's ability to produce
coherent predictions on real-world images.

\section{Analysis and Discussion}

\subsection{Innovation and Outreach Strategy}

EmotiAid represents a novel application of AI technology to a humanitarian context.
The key innovations are: (1) domain-specific deployment adapting a general FER
system to internally displaced children in Burkina Faso; (2) an interface
specifically designed for non-technical humanitarian agents; (3) an automatic
alert mechanism flagging high-risk emotional states. For outreach, EmotiAid could
be introduced to humanitarian organizations through partnerships with UNICEF and
the Burkina Faso Ministry of Social Action, with a pilot deployment in one IDP
camp for real-world validation.

\subsection{Limitations}

\begin{itemize}
    \item FER-2013 is primarily composed of adult faces; performance on child faces
    may differ from reported accuracy;
    \item The Disgust class shows lower performance due to severe dataset imbalance;
    \item The system requires a camera-equipped device and network connectivity;
    \item Cultural differences in emotional expression between Western training data
    and Sub-Saharan African subjects are not accounted for.
\end{itemize}

\subsection{Recommendations}

\begin{itemize}
    \item Implement Transfer Learning with MobileNetV2 to improve accuracy to
    72--75\%;
    \item Integrate AffectNet dataset for better ethnic diversity representation;
    \item Conduct a clinical validation study with child psychologists in a real
    IDP camp.
\end{itemize}

\subsection*{Chapter Conclusion}

This chapter presented the full implementation of EmotiAid. The FastAPI backend
successfully integrates the CNN model with real-time inference, JWT authentication,
and SQLite persistence. The React frontend provides a complete interface covering
analysis, history, child profiles, and reporting. The CNN achieved 67.53\%
validation accuracy on FER-2013, consistent with the academic state of the art.

\clearpage

%% ============================================================
%% GENERAL CONCLUSION
%% ============================================================
\chapter*{General Conclusion and Perspectives}
\addcontentsline{toc}{chapter}{General Conclusion and Perspectives}
\justifying

This thesis presented the design and implementation of EmotiAid, an AI-based web
application for the automatic detection of emotional states in internally displaced
children through facial expression analysis. The work was motivated by the critical
need for accessible psychosocial screening tools in humanitarian contexts,
particularly in Burkina Faso, where the ongoing security crisis has displaced more
than two million people.

A CNN trained on FER-2013 demonstrated the technical feasibility of automated
facial emotion classification with a validation accuracy of 67.53\%, exceeding the
academic standard of 65\% \cite{goodfellow2013}. The integration of this model
within a FastAPI/React web application produced a functional system enabling
non-specialist humanitarian agents to conduct rapid emotional assessments using a
standard mobile device.

The key contributions of this work are threefold. First, it demonstrates that AI
technology can be adapted and deployed for humanitarian applications in low-resource
environments. Second, it produces a functional prototype serving as a basis for
further development and clinical validation. Third, it articulates the technical
and ethical challenges that must be addressed for responsible deployment of such
systems with vulnerable children.

\textbf{Perspectives:} From a technical standpoint, the most immediate priority is
the integration of Transfer Learning with MobileNetV2 to improve model accuracy to
the 72--75\% range. From an organizational standpoint, a pilot deployment in a
Burkina Faso IDP camp in collaboration with a humanitarian organization and child
psychologists would constitute the critical next step toward real-world impact.

\clearpage

%% ============================================================
%% BIBLIOGRAPHY
%% ============================================================
\addcontentsline{toc}{chapter}{Bibliography}
\justifying

\begin{thebibliography}{99}

\bibitem{barsoum2016}
Barsoum E., Zhang C., Ferrer C.C. and Zhang Z. (2016),
Training deep networks for facial expression recognition with crowd-sourced label
distribution,
\textit{ACM International Conference on Multimodal Interaction}, pp. 279--283.

\bibitem{ekman1999}
Ekman P. (1999),
Basic emotions. In T. Dalgleish and M. Power (Eds.),
\textit{Handbook of Cognition and Emotion},
John Wiley \& Sons, pp. 45--60.

\bibitem{ekman1978}
Ekman P. and Friesen W.V. (1978),
\textit{Facial Action Coding System},
Consulting Psychologists Press, Palo Alto.

\bibitem{goodfellow2013}
Goodfellow I.J., Erhan D., Carrier P.L. and Courville A. (2013),
Challenges in representation learning: A report on three machine learning contests,
\textit{Neural Information Processing Systems Workshop}.

\bibitem{he2016}
He K., Zhang X., Ren S. and Sun J. (2016),
Deep residual learning for image recognition,
\textit{Proceedings of CVPR}, pp. 770--778.

\bibitem{lecun1998}
LeCun Y., Bottou L., Bengio Y. and Haffner P. (1998),
Gradient-based learning applied to document recognition,
\textit{Proceedings of the IEEE}, 86(11), pp. 2278--2324.

\bibitem{li2020}
Li S. and Deng W. (2020),
Deep facial expression recognition: A survey,
\textit{IEEE Transactions on Affective Computing}, 13(3), pp. 1195--1215.

\bibitem{mollahosseini2017}
Mollahosseini A., Hasani B. and Mahoor M.H. (2017),
AffectNet: A database for facial expression, valence, and arousal computing
in the wild,
\textit{IEEE Transactions on Affective Computing}, 10(1), pp. 18--31.

\bibitem{srivastava2014}
Srivastava N., Hinton G., Krizhevsky A., Sutskever I. and Salakhutdinov R. (2014),
Dropout: A simple way to prevent neural networks from overfitting,
\textit{Journal of Machine Learning Research}, 15(1), pp. 1929--1958.

\bibitem{taigman2014}
Taigman Y., Yang M., Ranzato M. and Wolf L. (2014),
DeepFace: Closing the gap to human-level performance in face verification,
\textit{Proceedings of CVPR}, pp. 1701--1708.

\bibitem{unhcr2023}
UNHCR (2023),
\textit{Global Trends: Forced Displacement in 2022},
United Nations High Commissioner for Refugees, Geneva.

\bibitem{viola2001}
Viola P. and Jones M. (2001),
Rapid object detection using a boosted cascade of simple features,
\textit{Proceedings of CVPR}, Vol. 1, pp. 511--518.

\bibitem{zhang2016}
Zhang K., Zhang Z., Li Z. and Qiao Y. (2016),
Joint face detection and alignment using multitask cascaded convolutional networks,
\textit{IEEE Signal Processing Letters}, 23(10), pp. 1499--1503.

\end{thebibliography}

\vspace{1cm}
\noindent\textbf{Websites}
\begin{enumerate}
    \item Kaggle FER-2013 Dataset. Available at:
    \url{https://www.kaggle.com/datasets/msambare/fer2013}. Accessed: March 2026.
    \item TensorFlow Documentation. Available at:
    \url{https://www.tensorflow.org/api_docs}. Accessed: March 2026.
    \item FastAPI Documentation. Available at:
    \url{https://fastapi.tiangolo.com}. Accessed: March 2026.
    \item React.js Documentation. Available at:
    \url{https://react.dev}. Accessed: March 2026.
    \item UNHCR Burkina Faso. Available at:
    \url{https://www.unhcr.org/countries/burkina-faso}. Accessed: March 2026.
\end{enumerate}

\clearpage

\end{document}